%% start of file `template.tex'.
%% Copyright 2006-2013 Xavier Danaux (xdanaux@gmail.com).
%
% This work may be distributed and/or modified under the
% conditions of the LaTeX Project Public License version 1.3c,
% available at http://www.latex-project.org/lppl/.

\documentclass[10pt,a4paper,sans]{moderncv} % possible options include font size ('10pt', '11pt' and '12pt'), paper size ('a4paper', 'letterpaper', 'a5paper', 'legalpaper', 'executivepaper' and 'landscape') and font family ('sans' and 'roman')

% moderncv themes
\moderncvstyle{classic} % style options are 'casual' (default), 'classic', 'oldstyle' and 'banking'
\moderncvcolor{orange} % color options 'blue' (default), 'orange', 'green',
\usepackage{fontawesome}
\providecommand\faSkype{{\FA\symbol{"F17E}}}

\usepackage[scale=0.80]{geometry}
\setlength{\hintscolumnwidth}{3cm} % if you want to change the column with the dates

\name{Giona}{Granchelli}
\title{Staff-Level Software Engineer | Java, Distributed Systems, Fintech}
\phone[mobile]{+31~(0)6~427433~16}
\email{giona.granchelli@gmail.com}
\social[linkedin]{giona-granchelli}
\social[github]{GionaGranchelli}
\website{\faGlobe\ https://gionag.com}
\photo[64pt][0.4pt]{face.jpg}
\quote{Chapter Lead and senior software engineer with 10+ years of experience designing and modernising distributed systems across banking, fintech, and cloud-native environments. Strong background in Java and Kotlin backend engineering, event-driven systems, Kafka, CQRS, data pipelines, and legacy-to-modern transformation. Hands-on technical leader focused on architecture direction, engineering enablement, technical debt reduction, and pragmatic adoption of AI-assisted development workflows to improve delivery quality and speed.}
%----------------------------------------------------------------------------------
% content
%----------------------------------------------------------------------------------
\begin{document}
\vspace*{-2cm}
\makecvtitle

\section{Core Skills}
\cvitem{Finance \& Modernisation}{
Legacy modernization, regulated systems, reconciliation, ETL, financial data processing, migration strategy
}
\cvitem{Backend Engineering}{
Java, Kotlin, Spring Boot, REST APIs, microservices, distributed systems
}
\cvitem{Data \& Integration}{
Kafka, CQRS, event-driven systems, schema evolution, data pipelines, asynchronous messaging
}
\cvitem{Cloud \& Platform}{
Kubernetes, Docker, Azure, AWS, CI/CD, observability, platform engineering
}
\cvitem{Frontend \& Product Delivery}{
Vue.js, React, TypeScript, Android, cross-functional product collaboration
}
\cvitem{Architecture \& Leadership}{
RFC-style technical communication, technical mentorship, engineering standards, system design trade-offs
}

\section{Selected Leadership Impact}
\cvitem{Technical Direction}{
Led modernization initiatives in banking and fintech environments, combining architecture work with hands-on implementation.
}
\cvitem{Cross-Team Enablement}{
Built reusable tooling, CI/CD patterns, and platform capabilities that improved deployment consistency and reduced delivery friction.
}
\cvitem{Distributed Systems}{
Applied event-driven patterns, Kafka-based integration, and CQRS approaches to support scalability, reliability, and clear service boundaries.
}
\cvitem{Mentorship}{
Supported engineers through code reviews, pairing, architectural discussions, and practical guidance on modern engineering practices.
}
\cvitem{AI-Assisted Engineering}{
Actively expanding AI-assisted and agentic development workflows for coding, testing, documentation, and developer productivity.
}

\section{Professional Experience}
\cventry{Oct 2024 - Now}{Chapter Lead Java \& VueJS | Software Engineer IV}
{ABN AMRO - BCDB}
{Amsterdam, Netherlands}
{}{
Leading technical modernization across a banking department responsible for legacy services and critical data flows, with a focus on cloud migration, engineering enablement, and reliability improvement.
\begin{itemize}
 \item Led modernization initiatives for legacy services, defining migration paths toward cloud-native architectures on Azure AKS.
 \item Designed Kubernetes-based platform patterns, CI/CD pipelines, and developer tooling to standardize containerization and deployment across teams.
 \item Built and improved reconciliation jobs validating financial data consistency across multiple banking systems in a regulated environment.
 \item Designed ETL pipelines and transformation flows for large-scale financial and master data processing used by downstream services and reporting.
 \item Improved reliability, observability, and operational robustness of batch and integration flows handling critical banking workloads.
 \item Acted as a technical multiplier for multiple teams by introducing reusable engineering patterns, paved-road style tooling, and migration guidance.
 \item Combined hands-on coding with technical leadership, supporting teams in modern engineering practices and reducing delivery friction during transformation.
 \item Stack: Kubernetes, Azure, AKS, Azure DevOps, Java, TypeScript, Kafka
\end{itemize}
}

\cventry{Sep 2020 - Sep 2024}{Full-stack | Software Engineer IV}
{ABN AMRO - Strategy and Innovation}
{Amsterdam, Netherlands}
{}{
Core engineer for PayDay, a fintech platform enabling instant payouts for gig-economy workers and freelancers.
\begin{itemize}
 \item Designed and implemented Kotlin and Spring Boot backend services for payment and payout workflows in a regulated financial environment.
 \item Built and maintained frontend and mobile capabilities supporting operational finance flows and end-user interactions.
 \item Contributed to architecture decisions around transaction handling, system integration, and scalable service boundaries.
 \item Delivered distributed services supporting financial transactions, workflow reliability, and user-facing product capabilities.
 \item Worked across backend, frontend, and mobile layers, combining product delivery with technical problem-solving in a high-trust domain.
 \item Stack: Kotlin, Spring Boot, Vue.js, Android, Docker, Kubernetes, AWS S3, PostgreSQL, GitLab CI
\end{itemize}
}

\cventry{May 2020 - Jul 2020}{Software Engineer}
{Ximedes}
{Haarlem, Netherlands}
{}{
Worked on a project for the Dutch public transportation sector (NS), contributing to an end-to-end solution enabling travel with QR codes instead of conventional tickets.
\begin{itemize}
 \item Built backend and integration components for ticketing workflows.
 \item Contributed to real-time communication features using WebSockets.
 \item Stack: Kotlin, Ktor, REST, WebSocket, ReactJS, TypeScript, Redis
\end{itemize}
}

\cventry{Oct 2019 - Apr 2020}{Full Stack Developer}
{BLOX - BTC Direct}
{Amsterdam / Nijmegen, Netherlands}
{}{
Worked on BLOX, a cryptocurrency trading platform serving approximately 500k daily users.
\begin{itemize}
 \item Maintained and extended microservices supporting trading, account management, and analytics.
 \item Contributed to a distributed event-driven architecture using Axon Framework and CQRS patterns.
 \item Improved stability and performance across multiple backend services.
 \item Stack: Java 8, Kotlin, Spring Boot, REST, gRPC, Axon Framework, Node.js, ReactJS, Cordova, MariaDB, Docker, Kubernetes, GitLab CI
\end{itemize}
}

\cventry{Jan 2019 - Sep 2019}{Senior Software Developer}
{WoodWing B.V.}
{Zaandam, Netherlands}
{}{
Worked on Elvis, a scalable distributed Digital Asset Management system for creating, managing, and distributing digital assets.
\begin{itemize}
 \item Developed and maintained backend and frontend features across a mature product environment.
 \item Participated in code reviews and peer programming, supporting engineering quality and shared ownership.
 \item Implemented an SSO solution integrating third-party and internal applications, including AWS Cognito and Okta.
 \item Contributed to the gradual split of a monolithic system by extracting authentication and authorization capabilities.
 \item Worked with a custom CQRS-style approach built around a modified use of Elasticsearch, supporting system evolution and read/write separation.
 \item Stack: Java 8, Spring Framework, Elasticsearch, AWS, Jenkins, Docker, REST, AngularJS, Node.js, Okta
\end{itemize}
}

\cventry{Jul 2017 - Jan 2019}{Full Stack Developer}
{Trifork B.V.}
{Amsterdam, Netherlands}
{}{
Worked on multiple projects in small agile consultancy teams, combining backend delivery, frontend work, and rapid onboarding into new domains.
\begin{itemize}
 \item Contributed to IBE, a cloud-based teacher-student platform for creating, managing, and taking school exams used in Switzerland.
 \item Developed new features and maintained legacy functionality in a product with active end users and ongoing delivery needs.
 \item Built the frontend application from scratch.
 \item Migrated and rebuilt CI/CD pipelines from Jenkins to GitLab.
 \item Worked on BLOX in a consultancy capacity before later joining directly.
 \item Developed an end-to-end service for managing supported coins on the platform.
 \item Built functionality for trade trend analysis and usage statistics.
 \item Stack: J2EE, Kotlin, Spring Boot, REST, gRPC, Axon Framework, Node.js, ReactJS, Webpack, Babel, Docker, Kubernetes, Jenkins, GitLab CI
\end{itemize}
}

\section{Education}
\cventry{Sep 2015 - Mar 2017}{Master Degree in Computer Science}
{University of L'Aquila}
{L'Aquila (AQ), Italy}
{maximum score, cum laude}{
Software Architecture, Embedded Systems, Model Driven Engineering, Research Methodologies, Formal Methods, and Mobile Applications.
\\ Thesis: Architecture Recovery of Microservice-based Systems.
\\ Supervisors: dr. Amleto di Salle, dr. Ivano Malavolta, and dr. Paolo di Francesco.
}

\cventry{Dec 2013 - Jun 2014}{First Level Postgraduate Master Degree in Web Technologies}
{University of L'Aquila}
{L'Aquila (AQ), Italy}
{}{
Completed coursework in LAMP, J2EE, UML for the web, ORM, XML, web services, SOA, web mining, and SEO.
}

\cventry{Sep 2009 - Jul 2013}{Bachelor Degree in Computer Science}
{University of L'Aquila}
{L'Aquila (AQ), Italy}
{99/110}{
Completed coursework in software engineering, Java and C programming, databases, web technologies, algorithms, compilers, and mathematics.
\\ Thesis: Automation control and reachable degree of automation through usage of commercial drones.
\\ Supervisors: dr. Davide Di Ruscio di Salle, dr. Ivano Malavolta, and dr. Patrizio Pellicione.
}

\section{Publications}
\cvitem{[1]}{MicroART: A Software Architecture Recovery Tool for Maintaining Microservice-based Systems. In Proceedings of the 14th International Conference on Software Architecture (ICSA). IEEE, 2017.}
\cvitem{[2]}{Towards Recovering the Software Architecture of Microservice-based Systems. In IEEE International Workshop on Architecting with MicroServices (AMS), April 2017.}

\section{Quote}
\cvitem{}{``...if you aren't, at any given time, scandalized by the code you wrote five or even three years ago, you're not learning anywhere near enough'' -- Nick Black}

\end{document}
